% Section 5, rewrite of October 21, 2026. Stationary solves at the estimated
% point (loss 7.4186) and the 2007 price 0.7707, rebate balanced in each solve.
% Figures and values: output/model/experiments/first_child_cost/20261021/
% (fig_s5_*.pdf, s5_values.json, scripts/s5_figs.py).
\newcommand{\FCC}{../../output/model/experiments/first_child_cost/20261021}

\section{Household Choices in the Calibrated Model}
\label{sec:results}

This section describes how households in the calibrated model respond to parenthood, and how earnings risk shapes the decision to have one. First, I show how having children changes housing, consumption and
tenure. The child's space trades off with consumption, so the cost
weighs most on households with little wealth or low earnings. Second, I illustrate how
 lowering the standard deviation of earnings shocks by about a quarter,
holding the house price at its 2007 level, raises completed fertility by 12\%.

\subsection{Choices at the first birth}
\label{sec:results-policy-functions}

Figure~\ref{fig:results-policy} compares the choices of the same household in
the period in which it has no child and in the period in which a first child
is born. The household is a renter aged 26 to 29 with median earnings, and
the figure varies its liquid wealth. Without savings, a first child pushes the
space the household rents close to the rental cap and lowers the probability
of owning from 42\% to 24\%. The child needs space when the household can
least afford a down payment, so it rents and pays for the space out of
consumption. With more than about two years of earnings in liquid wealth, the
child instead raises the probability of owning, and the household buys a
larger home.

Panel (d) shows the value of trying for a child, $V^F_{t,A}-V^F_{t,0}$. It is
negative without savings, turns positive at about five months of earnings,
and rises steeply at low wealth before flattening. Wealth therefore matters
for the decision to have a first child mostly among households with little
of it.

\begin{figure}[!tb]
\centering
\includegraphics[width=0.72\linewidth]{\FCC/fig_s5_choices.pdf}
\caption{Choices with and without a first child}
\label{fig:results-policy}
\par\smallskip
\begin{minipage}{0.92\linewidth}
\footnotesize\textit{Notes:} Renter aged 26--29 with median earnings
($z=0.78$), 2007 stationary equilibrium. Liquid wealth in units of mean annual earnings.
Rooms and consumption are averages over tenure choices.
\end{minipage}
\end{figure}

\subsection{Choices across earnings}
\label{sec:results-cost}

Figure~\ref{fig:results-cost} repeats the comparison across earnings states,
for renters aged 26 to 29 without savings. A first child adds about two thirds
of a room at earnings up to the median, and about half a room at high earnings,
where most of these households buy a home even without a child. Consumption falls
by 21\% at the lowest state and by 2\% at the median state. At high earnings
it rises, because the equivalence scale raises the household's needs and
these households can afford to meet them. As Section~\ref{sec:m-analysis} shows, the space requirement alone
costs about $r_th_P$ in every period the child is at home, whether the
household rents or owns. At the 2007 price, this is 2.9\% of mean annual
earnings each year. The cost of a child's space is therefore a recurring
expense, and it weighs most on households with low earnings.

\begin{figure}[!b]
\centering
\includegraphics[width=\linewidth]{\FCC/fig_s5_cost_by_earnings_all9.pdf}
\caption{Choices with and without a first child, by earnings}
\label{fig:results-cost}
\par\smallskip
\begin{minipage}{0.92\linewidth}
\footnotesize\textit{Notes:} Renters aged 26--29 without savings, 2007
stationary equilibrium, by earnings state $z$. Rooms and consumption are
averages over tenure choices. The dashed line marks the rental cap $\bar h^R$.
\end{minipage}
\end{figure}

\subsection{Earnings risk}
\label{sec:results-risk}

Section~\ref{sec:m-analysis} showed that earnings risk raises the expected
cost of a child, since that cost falls with earnings at a decreasing rate. To
measure this force, I lower the standard deviation of the four-year earnings
shock by about a quarter, from 0.48 to 0.35, rescale the earnings states to keep mean earnings
unchanged, and solve the model at the 2007 price. Lower risk raises completed
fertility by 12\%, from 2.10 to 2.36.\footnote{Without the rescaling, the same
change raises mean earnings by 4.9\%, and completed fertility rises by 16\%.}
Figure~\ref{fig:results-risk-floor} below reports the effect over a range of
values of earnings risk.

Figure~\ref{fig:results-risk} shows where the additional births come from.
In the baseline, the model matches the distribution of children ever born
among women in their early forties. Lower risk shifts households out of
childlessness and toward larger families: the share without children falls
from 20\% to 12\%, and the shares with two and with three or more rise.

\begin{figure}[!tb]
\centering
\includegraphics[width=0.6\linewidth]{\FCC/fig_s5_risk_children.pdf}
\caption{Earnings risk and children ever born}
\label{fig:results-risk}
\par\smallskip
\begin{minipage}{0.92\linewidth}
\footnotesize\textit{Notes:} Shares with 0, 1, 2, and 3 or more children
ever born. Data: CPS June Fertility Supplement, 2006 and 2008, women aged
40--44. Model: households at about age 42, 2007 stationary equilibrium, at
the estimated earnings risk ($\sigma_z=0.48$, baseline) and at lower risk
($\sigma_z=0.35$, earnings states rescaled to keep mean earnings unchanged),
same price.
\end{minipage}
\end{figure}

Earnings risk weighs more when a child requires more space.
Figure~\ref{fig:results-risk-floor} solves the model at the 2007 price for
different values of the space a child requires, $h_P$, and of the standard
deviation of earnings shocks. Lowering risk to 0.35 raises completed
fertility by 10\% when a child requires half a room, and by 17\% when it
requires 1.6 rooms. The child's space must be paid in every period the child
is at home, including after a bad earnings draw, when consumption is most
valuable. The more space a child requires, the more these periods cost, and
the more a household gains from lower risk.

\begin{figure}[H]
\centering
\includegraphics[width=0.65\linewidth]{\FCC/fig_s5_risk_floor25.pdf}
\caption{Earnings risk and the space a child requires}
\label{fig:results-risk-floor}
\par\smallskip
\begin{minipage}{0.92\linewidth}
\footnotesize\textit{Notes:} Completed fertility, \% change from the
baseline (marked), 2007 stationary equilibrium, same price, on a grid of 25
values of each parameter. Vertical axis: standard deviation of the four-year
earnings shock, with earnings states rescaled to keep mean earnings
unchanged. Horizontal axis: space a child requires, $h_P$, in rooms.
\end{minipage}
\end{figure}
